% รายงานการประชุม — รูปแบบอ้างอิง กบ.วช. ครั้งที่ 8/2568 (145731สรุปบันทึกการประชุม...)
% หลักสูตร วศ.ม. วิศวกรรมคอมพิวเตอร์และปัญญาประดิษฐ์ ม.ร.อ.
\documentclass{meetingminutes}

\meetingcommittee{คณะกรรมการบริหารหลักสูตร วิศวกรรมคอมพิวเตอร์และปัญญาประดิษฐ์ (บัณฑิตศึกษา)}
\meetingnumber{๕/ภาคเรียนที่ ๒/๒๕๖๘}{๒๕๖๘}
\meetingdate{วันพฤหัสบดีที่ ๒๗ มีนาคม พ.ศ. ๒๕๖๙}
\meetingplace{ห้องประชุมคณะเทคโนโลยีอุตสาหกรรม}
\meetingstart{๑๓.๐๐ น.}
\meetingend{๑๖.๑๐ น.}

\begin{document}
\maketitle

\psection{ผู้มาประชุม}
\begin{participants}
  \pmember{1}{ผู้ช่วยศาสตราจารย์ ดร.พัชรี มณีรัตน์}{ประธานหลักสูตร}{ประธาน}
  \pmember{2}{ผู้ช่วยศาสตราจารย์ ดร.พิศิษฐ์ นาคใจ}{อาจารย์ประจำหลักสูตร}{กรรมการ}
  \pmember{3}{ผู้ช่วยศาสตราจารย์ ดร.กาญจนา ดาวเด่น}{อาจารย์ประจำหลักสูตร (เลขานุการ)}{กรรมการ/เลขานุการ}
\end{participants}

\psection{ผู้ไม่มาประชุม}
\noindent ไม่มี\par

\startmeeting

\openingchair{ผู้ช่วยศาสตราจารย์ ดร.พัชรี มณีรัตน์ ประธานการประชุม}{กล่าวเปิดการประชุมและดำเนินการตามระเบียบวาระ ดังนี้}

\agenda{๑}{เรื่องแจ้งเพื่อทราบ — รับรองรายงานการประชุมครั้งก่อน}

\subagenda{๑.๑}{รับรองรายงานการประชุมครั้งที่ ๔/ภาคเรียนที่ ๒/๒๕๖๘}

ที่ประชุมรับรองรายงานการประชุมครั้งที่ ๔/ภาคเรียนที่ ๒/ปีการศึกษา ๒๕๖๘ โดยไม่มีข้อแก้ไข

\resolution{รับรอง}

\agenda{๒}{เรื่องเสนอเพื่อพิจารณา — สรุปผลการจัดการศึกษาภาคเรียนที่ ๒/๒๕๖๘}

\subagenda{๒.๑}{สรุปผลการเรียนการสอนและการประเมิน CLO รายวิชาบังคับ}

อาจารย์ผู้สอนรายงานสรุปผลการจัดการศึกษาภาคเรียนที่ ๒/๒๕๖๘ โดยมีนักศึกษารุ่นแรกจำนวน ๔ คน ผลสรุปดังนี้

\begin{itemize}
  \item รายวิชา ๔๐๙๕๑๐๑ ปัญญาประดิษฐ์และการเรียนรู้ของเครื่อง — นักศึกษาทั้ง ๔ คนผ่านการประเมิน CLO ครบทุกข้อ
  \item รายวิชา ๔๐๙๕๑๐๒ การเขียนโปรแกรมขั้นสูงสำหรับการเรียนรู้ของเครื่อง — นักศึกษาทั้ง ๔ คนผ่านการประเมิน CLO ครบทุกข้อ รวมถึง CLO ด้านจริยธรรมในการพัฒนาโปรแกรม
  \item อัตราผ่านภาคเรียน (Pass Rate) เท่ากับ ๑๐๐ เปอร์เซ็นต์ ไม่มีนักศึกษาตกหรือลาออกในภาคเรียนนี้
\end{itemize}

\resolution{รับทราบ}

\subagenda{๒.๒}{รายงานกิจกรรม IEL Community Workshop วันที่ ๒๖ มีนาคม พ.ศ. ๒๕๖๙}

ผู้ช่วยศาสตราจารย์ ดร.พิศิษฐ์ นาคใจ รายงานว่านักศึกษาทุกคนมีส่วนร่วมถ่ายทอดความรู้ด้านปัญญาประดิษฐ์และเทคโนโลยีให้โรงเรียนในจังหวัดอุตรดิตถ์ จำนวน ๓ แห่ง สะท้อนผลการเรียนรู้ตาม PLO ๒ และ PLO ๕

\resolution{รับทราบ}

\subagenda{๒.๓}{รายงานความก้าวหน้าหัวข้อวิจัย/สารนิพนธ์ และการมอบหมายอาจารย์ที่ปรึกษา}

ที่ประชุมรับทราบว่านักศึกษาทุกคนมีหัวข้อวิจัย/สารนิพนธ์ชัดเจนแล้ว โดยสรุปดังนี้

\begin{longtable}{|p{3.2cm}|p{5.5cm}|p{4.8cm}|}
  \hline
  \textbf{นักศึกษา} & \textbf{หัวข้อวิจัย/สารนิพนธ์} & \textbf{อาจารย์ที่ปรึกษา} \\
  \hline
  นางสาวพอใจ สุขหา & IoT ตู้อบแห้งข้าวแต๋นอัจฉริยะ & ผศ.ดร.พัชรี มณีรัตน์ \\
  \hline
  นายจิรกิตติ์ วราภรณ์ & Remote Sensing ระบุพื้นที่ปลูกอ้อย & ผศ.ดร.พิศิษฐ์ นาคใจ \\
  \hline
  นายจาตุรงค์ ทรงภักดี & AI จัดการสถานีชาร์จ EV & อาจารย์ที่ปรึกษาภายนอก \\
  \hline
  นายเกริกเกียรติ บุญไทย & Deep Learning ตรวจจับผู้ไม่สวมหมวกนิรภัย & ผศ.ดร.กาญจนา ดาวเด่น \\
  \hline
\end{longtable}

\resolution{รับทราบหัวข้อวิจัยของนักศึกษาทุกราย และยืนยันการมอบหมายอาจารย์ที่ปรึกษาตามตาราง โดยมอบหมายให้อาจารย์ที่ปรึกษานัดพบนักศึกษาเพื่อวางแผนความก้าวหน้าในภาคเรียนที่ ๓/๒๕๖๘}

\agenda{๓}{เรื่องเสนอเพื่อพิจารณา — วางแผนภาคเรียนที่ ๓/๒๕๖๘ และการรับสมัครปีการศึกษา ๒๕๖๙}

\subagenda{๓.๑}{แผนการจัดการศึกษาภาคเรียนที่ ๓/๒๕๖๘ (ภาคฤดูร้อน)}

ที่ประชุมวางแผนเน้นความก้าวหน้าวิทยานิพนธ์/สารนิพนธ์ของนักศึกษาทั้ง ๔ คน ตามรายวิชาที่ลงทะเบียน

\resolution{เห็นชอบ}

\subagenda{๓.๒}{เตรียมเปิดรับสมัครนักศึกษารุ่นที่ ๒ ปีการศึกษา ๒๕๖๙}

ประธานแจ้งว่าจะดำเนินการเปิดรับสมัครนักศึกษาในช่วงภาคเรียนที่ ๓/๒๕๖๘ และกำหนดประชุมติดตามในเดือนเมษายน พ.ศ. ๒๕๖๙

\resolution{เห็นชอบทุกวาระ และกำหนดประชุมครั้งต่อไปในเดือนเมษายน พ.ศ. ๒๕๖๙}

\adjournmeeting

\begin{sigblock}
  \typist{ผู้ช่วยศาสตราจารย์ ดร.กาญจนา ดาวเด่น}
  \reviewer{ผู้ช่วยศาสตราจารย์ ดร.พิศิษฐ์ นาคใจ}{กรรมการ}
  \chairsig{ผู้ช่วยศาสตราจารย์ ดร.พัชรี มณีรัตน์}{ประธานการประชุม}
\end{sigblock}

\end{document}
